\section{多头注意力（Multi-Head Attention）}

\subsection{动机与设计}

单个 Self Attention 层使用一组 $W_Q, W_K, W_V$ 矩阵。为了增加模型的容量和表达能力，我们可以运行 $H$ 个\textbf{独立}的 Self Attention``头''（head）。

\begin{figure}[H]
\centering
\includegraphics[width=0.95\textwidth]{slides-images/slide-039.jpg}
\caption{Multi-Head Attention：$H$ 个独立的 Self Attention 头并行处理相同的输入，各自有不同的权重矩阵，输出拼接后通过线性投影融合\protect\footnotemark}
\end{figure}
\footnotetext{Slides 第40页。}

\begin{importantbox}{Multi-Head Attention 的计算}
\begin{enumerate}
    \item 对输入 $X$，每个头 $h$ 独立计算 $Q^h, K^h, V^h$（使用各自的权重矩阵）
    \item 每个头独立执行注意力：$Y^h = \text{Attention}(Q^h, K^h, V^h)$
    \item 将所有头的输出拼接：$[Y^1; Y^2; \ldots; Y^H]$
    \item 通过一个线性投影 $W_O$ 融合所有头的信息
\end{enumerate}

每个头可以学习关注输入的不同方面：
\begin{itemize}
    \item 某个头可能学会关注语法结构
    \item 另一个头可能关注语义相似性
    \item 还有的头可能关注局部上下文
\end{itemize}
\end{importantbox}

\subsection{计算效率}

尽管 Multi-Head Attention 看起来很复杂，但整个操作本质上只是\textbf{四次（批量）矩阵乘法}：

\begin{enumerate}
    \item \textbf{QKV 投影}：一次矩阵乘法，将输入投影为所有头的 Q、K、V
    \item \textbf{Q-K 相似度}：一次批量矩阵乘法，计算所有头的注意力分数
    \item \textbf{V 加权}：一次批量矩阵乘法，对 Value 进行加权求和
    \item \textbf{输出投影}：一次矩阵乘法，融合所有头的输出
\end{enumerate}

\begin{knowledgebox}{实际实现中的高效技巧}
在实际代码中，$H$ 个头的计算并不需要 for 循环，而是通过\textbf{张量重塑}（reshape）和\textbf{批量矩阵乘法}（batched matmul）来并行计算。三个投影矩阵 $W_Q, W_K, W_V$ 也通常拼接为一个大矩阵，用一次矩阵乘法完成所有投影。这使得 Multi-Head Attention 在 GPU 上非常高效。
\end{knowledgebox}

\subsection{本章小结}

Multi-Head Attention 通过运行多个独立的注意力头来增加模型容量。每个头学习不同的注意力模式，最终通过线性投影融合。整个操作高度适合 GPU 并行计算。

%% ========== 三大序列处理原语对比 ==========

